% Source: https://artofproblemsolving.com/community/c3h2071105_roman_numerals__2019_purple_comet_10
% Stable identity: aops-topic-2071105-q01
The student is wrong because Roman numerals are not just sums: in the subtractive pairs $IV,IX,XL,XC,CD,CM$ the smaller symbol written first is \emph{subtracted}. Rearranging symbols can therefore change the value; for example $VI=6$ but $IV=4$.

\emph{Greatest.} Using every symbol additively gives $MDCLXVI=1000+500+100+50+10+5+1=1666$, the largest possible since each symbol contributes at most its face value. So $N=1666$.

\emph{Least.} Only $I$, $X$ and $C$ can ever be subtracted (and only in the listed pairs). $M$ can never be subtracted, so it contributes $1000$. Of $D$ and $C$, the only way to use both with $C$ subtracted is $CD=400$ ($CM$ is excluded because $D$ appears); otherwise they give $600$. Similarly $L$ and $X$ give at least $XL=40$, and $V$ and $I$ give at least $IV=4$. The numeral $MCDXLIV$ is valid and achieves all minima, so $n=1000+400+40+4=1444$.

Therefore $N-n=1666-1444=\boxed{222}$.
